\chapter{Conclusiones}

La extensión del esquema de emisión de paquetes de $n$ fonones al caso de una molécula excitónica demostró que el acoplamiento colectivo introducido por la interacción de Förster modifica cualitativamente la dinámica del sistema respecto al caso de un único emisor. El factor superradiante $\sqrt{2}$ y el corrimiento $-J$ en las resonancias de Stokes no son correcciones perturbativas menores sino efectos estructurales que determinan tanto la posición de las resonancias como la escala temporal de la cascada de emisión. Esta distinción es la clave para entender por qué la molécula excitónica puede operar con alta pureza bajo condiciones experimentales menos exigentes.

La pureza de emisión $\Pi_3 > 95\%$ es accesible en la molécula excitónica con factores de calidad $Q \sim 90$--$300$, frente a $Q \sim 500$ requeridos por el caso de un único punto cuántico. Esta ventaja no proviene de un ajuste fino de parámetros sino del mecanismo físico subyacente: al amplificar $\Omega_\mathrm{eff}^{(n)}$ por el factor $\sqrt{2}$, la cascada de emisión ocurre más rápidamente y la ventana temporal disponible para que los procesos fuera de resonancia contaminen el paquete se reduce. El sistema es entonces inherentemente más robusto frente a las pérdidas de la cavidad.

La estadística de los paquetes emitidos es sintonizable entre tres regímenes físicamente distintos n-phonon gun, n-phonon laser y estado térmico de paquetes, mediante el ajuste de la relación $\gamma/\kappa$. Esta sintonizabilidad, combinada con la posibilidad de detectar el fotón heraldo que anuncia cada evento de emisión, hace del esquema una plataforma flexible para distintos protocolos de información cuántica fonónica.

Finalmente, los parámetros empleados son consistentes con los reportados experimentalmente en sistemas de puntos cuánticos acoplados a nanocavidades acústicas de THz, lo que sitúa las predicciones de este trabajo dentro del alcance de la tecnología actual y abre una ruta concreta hacia la realización experimental de fuentes acústicas cuánticas multifonónicas colectivas. En esta dirección, como perspectivas para trabajo futuro, resulta de interés extender el análisis a sistemas de $N > 2$ puntos cuánticos, donde los efectos colectivos pueden amplificarse progresivamente con $N$, y estudiar el impacto de la asimetría entre los emisores sobre la dinámica del estado oscuro. La incorporación de desfase puro mediante trayectorias cuánticas exactas de Monte Carlo permitiría además superar la limitación metodológica del método de Muñoz identificada en este trabajo para el cálculo de $\Pi_n$ en presencia de $\gamma_\varphi$.